
\section{弱相互作用玻色费米气体}

\subsection{硬球赝势的引入}

对于二体碰撞，取相对坐标 \(\vb{r} = \vb{r}_{i} - \vb{r}_{j}\)，记折合质量为\(M\)，硬球势参量为\(a\)，则波函数满足
\begin{equation*}
	\left( -\frac{\hbar^{2}}{2M} \nabla^{2} + V(r) \right) \psi(\vb{r}) = E \psi(\vb{r}), \quad E = \frac{\hbar^{2} k^{2}}{2M}, \quad V(r) =
	\begin{cases}
		+\infty, & r < a, \\
		0,       & r > a.
	\end{cases}
\end{equation*}

这里 \(a\) 是硬球半径；在最低阶 \(s\) 波近似中，它同时充当散射长度。零程近似要求相对动量满足 \(ka\ll1\)，多体热力学中还要满足 \(a/\lambda\ll1\)。

显然，在硬球内部 (\(r < a\))，\(\psi(\vb{r}) = 0\)，在硬球外部 (\(r > a\))，解为
\begin{equation*}
	\psi(\vb{r}) = \sum_{l=0}^{\infty} \sum_{m=-l}^{l} \psi_{lm}(\vb{r}), \quad \psi_{lm}(\vb{r}) = C_{lm} \left( j_{l}(kr) - \tan\delta_{l} \, n_{l}(kr) \right) Y_{lm}(\uv{r}).
\end{equation*}
式中\(C_{lm}\) 是常数，相移量为
\begin{equation*}
	\delta_{l} = \arctan\left( \frac{j_{l}(ka)}{n_{l}(ka)} \right).
\end{equation*}

如上外部解对于\(r > 0\) 的自由空间显然成立，而在原点波函数发散
\begin{equation*}
	\psi_{lm}\big|_{r \to 0} = \frac{C_{lm} k^{l}}{(2l+1)!!} \left( r^{l} + \frac{(2l+1)!!(2l-1)!!}{k^{2l+1}} \tan\delta_{l} \, r^{-l-1} \right) Y_{lm}(\uv{r}).
\end{equation*}

因此，若想将外部解延拓至\(r = 0\)，须在原点引入含\(\delta\)函数发散项的赝势。具体而言，对于原点附近无穷小体积的积分，不难证明
\begin{equation*}
	\int_{V_{\varepsilon}} r^{l} \left( \frac{1}{r^{2}} \frac{\d}{\d r} \left( r^{2} \frac{\d}{\d r} \right) - \frac{l(l+1)}{r^{2}} \right) n_{l}(kr) \, \d \vb{r} = \int_{V_{\varepsilon}} \left( r^{l} \nabla^{2} n_{l}(kr) - n_{l}(kr) \nabla^{2} r^{l} \right) \d \vb{r} = \frac{4\pi}{k^{l+1}} (2l+1)!!.
\end{equation*}
给出
\begin{equation*}
	\left( \nabla^{2} + k^{2} \right) \left( n_{l}(kr) Y_{lm}(\uv{r}) \right) = \frac{4\pi}{k^{l+1} r^{l}} (2l+1)!! \, \delta^3(\vb{r}) Y_{lm}(\uv{r}).
\end{equation*}

进而，薛定谔方程中发散部分为
\begin{equation*}
	\left( \nabla^{2} + k^{2} \right) \psi_{lm} \big|_{\mathrm{div}} = -C_{lm} \tan\delta_{l} \left( \nabla^{2} + k^{2} \right) \left( n_{l}(kr) Y_{lm}(\uv{r}) \right) = -\frac{4\pi (2l+1)!!}{k^{l+1} r^{l}} \tan\delta_{l}\, C_{lm} Y_{lm}(\uv{r}) \delta^3(\vb{r}).
\end{equation*}

代入
\begin{equation*}
	C_{lm} Y_{lm}(\uv{r}) = \frac{1}{(2l)!! k^{l}} \left( \frac{\partial^{2l+1}}{\partial r^{2l+1}} \left( r^{l+1} \psi_{lm}(\vb{r}) \right) \right) \bigg|_{r \to 0}.
\end{equation*}
得到赝势为
\begin{equation*}
	V_{\mathrm{ps}} = -\frac{2\pi \hbar^{2}}{M} \sum_{l=0}^{\infty} \frac{(2l+1)!!}{(2l)!!} \delta^3(\vb{r}) \frac{\tan\delta_{l}}{k^{2l+1}} \left( \frac{\partial^{2l+1}}{\partial r^{2l+1}} r^{l+1} \mathcal{P}_{l}(\uv{r}) \right).
\end{equation*}
其中 \(\mathcal{P}_{l}\) 是投影算符
\begin{equation*}
	\mathcal{P}_{l}(\uv{r}) = \sum_{m=-l}^{l} Y_{lm}(\uv{r}) \int \psi(\vb{r}') Y_{lm}^{*}(\uv{r}') \, \d \Omega'.
\end{equation*}
这里采用复球谐函数的通常归一化，投影系数必须含 \(Y_{lm}^{*}\)；若改用实球谐基，则复共轭退化为相同的实函数.

对于低能量散射 (\(ka \ll 1\))
\begin{equation*}
	-\frac{\tan\delta_{l}}{k^{2l+1}} \bigg|_{ka \ll 1} \simeq \frac{a^{2l+1}}{(2l+1)!!(2l-1)!!}.
\end{equation*}
赝势展开为
\begin{equation*}
	V_{\mathrm{ps}} &= \frac{2\pi \hbar^{2}}{M} \sum_{l=0}^{\infty} \frac{a^{2l+1}}{(2l)!!(2l-1)!!} \delta^3(\vb{r}) \left( \frac{\partial^{2l+1}}{\partial r^{2l+1}} r^{l+1} \mathcal{P}_{l}(\uv{r}) \right)\\&= \frac{2\pi \hbar^{2} a}{M} \delta^3(\vb{r}) \left( \frac{\partial}{\partial r} r + \frac{a^{2}}{2} \frac{\partial^{3}}{\partial r^{3}} r^{2} \mathcal{P}_{1}(\uv{r}) + \frac{a^{4}}{24} \frac{\partial^{5}}{\partial r^{5}} r^{3} \mathcal{P}_{2}(\uv{r}) + \mathcal{O}(a^{6}) \right).
\end{equation*}

通常，对于二体碰撞，记单个粒子质量为\(m\)，则\(M = m/2\)，保留到硬球势参量的最低阶项，有
\begin{equation*}
	V_{\mathrm{ps}} = \frac{4\pi \hbar^{2} a}{m} \delta^3(\vb{r}) \left( \frac{\partial}{\partial r} r \right).
\end{equation*}

下面多体计算只保留 \(l=0\) 接触项，适用条件仍为 \(a/\lambda\ll1\)。

进而哈密顿量修正为
\begin{equation*}
	\mathcal{H} = \mathcal{H}_{0} + \mathcal{H}_{1}, \quad \mathcal{H}_{0} = -\frac{\hbar^{2}}{2m} \sum_{i} \nabla_{i}^{2}, \quad \mathcal{H}_{1} = \frac{4\pi a \hbar^{2}}{m} \sum_{i<j} \delta(\vb{r}_{i} - \vb{r}_{j}) \left( \frac{\partial}{\partial r_{ij}} r_{ij} \right) = \frac{2\pi a \hbar^{2}}{m} \sum_{i\ne j} \delta(\vb{r}_{i} - \vb{r}_{j}) \left( \frac{\partial}{\partial r_{ij}} r_{ij} \right).
\end{equation*}
其中\(i, j\)为粒子编号， \(r_{ij} = |\vb{r}_{i} - \vb{r}_{j}|\), \(\nabla_{i}^{2}\)以 \(\vb{r}_{i}\) 为微分坐标. 第二种写法是把无序粒子对改写成有序但互异的粒子对，因此必须排除 \(i=j\) 的自相互作用.

\subsection{理想玻色费米气体：统计函数与低温基准}

先暂时令 \(\mathcal{H}_{1}=0\)，建立后文所有相互作用修正都要用到的理想气体基准。玻色子和费米子的差异在这里集中表现为占据数的取值范围以及相应的统计函数；自旋简并度则作为态密度的乘子保留。为了同时覆盖非相对论、相对论和一般幂律色散，以下先使用累计态数 \(\Sigma(\varepsilon)\) 的一般形式.

\subsubsection{巨配分函数与统一统计函数}

理想玻色费米气体的巨配分函数为
\begin{equation*}
	\Xi = \sum_{N=0}^{\infty} z^{N} \sum_{r} e^{-\beta E_{r}} = \prod_{\alpha} \sum_{n_{\alpha}} \left( z e^{-\beta \varepsilon_{\alpha}} \right)^{n_{\alpha}}.
\end{equation*}
其中，\(r\) 代表 \(\{n_{\alpha}\}\) 分布的序号数，粒子数满足约束 \(N = \sum_{\alpha} n_{\alpha}\)，能量可表示为 \(E = \sum_{\alpha} n_{\alpha} \varepsilon_{\alpha}\).

对于玻色子，\(n_{\alpha}\) 可以取一切自然数，对于费米子，\(n_{\alpha} = 0, 1\)，故
\begin{equation*}
	\Xi =
	\begin{cases}
		\prod_{\alpha} \left( 1 - z e^{-\beta \varepsilon_{\alpha}} \right)^{-1}, & \text{B.E.} \\
		\prod_{\alpha} \left( 1 + z e^{-\beta \varepsilon_{\alpha}} \right),      & \text{F.D.}
	\end{cases}
	\quad
	\ln \Xi =
	\begin{cases}
		-\sum_{\alpha} \ln \left( 1 - z e^{-\beta \varepsilon_{\alpha}} \right), & \text{B.E.} \\
		\sum_{\alpha} \ln \left( 1 + z e^{-\beta \varepsilon_{\alpha}} \right),  & \text{F.D.}
	\end{cases}
\end{equation*}

在热力学极限下，可以将\(d\)维空间的求和连续化为\(2d\)维相空间的积分
\begin{equation*}
	\varepsilon_{\alpha} \to \varepsilon, \quad \sum_{\alpha} \to \frac{1}{h^{d}} \int \prod_{i=1}^{d} \d q_{i} \, \d p_{i}.
\end{equation*}

\subsubsection{一般态密度与热力学闭式公式}

对于自由粒子系统，以上对位置空间的积分得到体积\(V\)，对动量空间的积分可以转换至能量空间，如上积分总可以写作
\begin{equation*}
	\frac{1}{h^{d}} \int \prod_{i=1}^{d} \d q_{i} \, \d p_{i} = \int \frac{\d \Sigma(\varepsilon)}{\d \varepsilon} \, \d \varepsilon.
\end{equation*}

不妨假设 \(\Sigma(\varepsilon) = B V \varepsilon^{s}, s \ge 1\)，\(B\) 是与体积无关的常数，进而
\begin{equation*}
	\sum_{\alpha} \to B V \int \d \varepsilon^{s} = B V \int_{0}^{\infty} s \varepsilon^{s-1} \, \d \varepsilon.
\end{equation*}

\begin{notes}
	比如，对于\(d\)维空间中能谱形如 \(\varepsilon = \kappa p^{a}\) 且内禀简并度为\(g\)的粒子
	\begin{equation*}
		\Sigma(\varepsilon) = \frac{g V}{h^{d}} \frac{\Omega_{d}}{d} p^{d} = g \frac{\Omega_{d}}{d} \frac{V}{h^{d}} \left( \frac{\varepsilon}{\kappa} \right)^{d/a}, \quad \Omega_{d} = \frac{2\pi^{d/2}}{\Gamma(d/2)}.
	\end{equation*}
	给出
	\begin{equation*}
		B=\frac{g}{d}\frac{\Omega_d}{h^d}\kappa^{-d/a},\quad s=\frac{d}{a}
	\end{equation*}
	质量为\(m\)的经典粒子对应 \(a = 2, \kappa = (2m)^{-1}\)，极端相对论型粒子对应 \(a = 1, \kappa = c\)。在一些复杂结构中，\(\kappa, \alpha\) 会有更为特殊的取值。
\end{notes}

为统一表示玻色和费米统计，记
\[
\Phi_{\nu}(z)=
\begin{cases}
g_{\nu}(z), & \text{B.E.},\\
f_{\nu}(z), & \text{F.D.}
\end{cases}
\]
其中 \(g_{\nu}\) 和 \(f_{\nu}\) 的定义域分别由玻色逸度 \(0\le z\le1\) 与费米逸度 \(z>0\) 决定。代入巨配分函数，转换求和为积分得到
\begin{equation*}
	\ln \Xi &= B V \Gamma(s+1) (kT)^{s} \Phi_{s+1}(z). \quad
	N = z \frac{\partial \ln \Xi}{\partial z} = B V \Gamma(s+1) (kT)^{s} \Phi_{s}(z). \\
	E &= -\frac{\partial \ln \Xi}{\partial \beta} = s N k T \frac{\Phi_{s+1}(z)}{\Phi_{s}(z)}. \quad
	P V = k T \ln \Xi = N k T \frac{\Phi_{s+1}(z)}{\Phi_{s}(z)}. \\
	S &= k \left( \ln \Xi - \beta \frac{\partial \ln \Xi}{\partial \beta} - z \ln z \frac{\partial \ln \Xi}{\partial z} \right) = N k \left( -\ln z + (s+1) \frac{\Phi_{s+1}(z)}{\Phi_{s}(z)} \right).
\end{equation*}

沿着等体积线
\begin{equation*}
	&\left( \frac{\partial z}{\partial T} \right)_{V} = -\frac{s z}{T} \frac{\Phi_{s}(z)}{\Phi_{s-1}(z)}. \quad C_{V} = T \left( \frac{\partial S}{\partial T} \right)_{V} = s N k \left( (s+1) \frac{\Phi_{s+1}(z)}{\Phi_{s}(z)} - s \frac{\Phi_{s}(z)}{\Phi_{s-1}(z)} \right). \\
	&\left( \frac{\partial C_{V}}{\partial T} \right)_{V} = \frac{s^{2} N k}{T} \left( (s+1) \frac{\Phi_{s+1}(z)}{\Phi_{s}(z)} - \frac{\Phi_{s}(z)}{\Phi_{s-1}(z)} - s \frac{\Phi_{s}^{2}(z) \Phi_{s-2}(z)}{\Phi_{s-1}^{3}(z)} \right).
\end{equation*}

沿着等压线
\begin{equation*}
	\left( \frac{\partial z}{\partial T} \right)_{P} = -\frac{(s+1) z}{T} \frac{\Phi_{s+1}(z)}{\Phi_{s}(z)}. \quad C_{P} = T \left( \frac{\partial S}{\partial T} \right)_{P} = (s+1) N k \left( (s+1) \frac{\Phi_{s-1}(z) \Phi_{s+1}^{2}(z)}{\Phi_{s}^{3}(z)} - s \frac{\Phi_{s+1}(z)}{\Phi_{s}(z)} \right).
\end{equation*}

沿着等温线
\begin{equation*}
	\left( \frac{\partial z}{\partial V} \right)_{T} = -\frac{z}{V} \frac{\Phi_{s}(z)}{\Phi_{s-1}(z)}. \quad
	\kappa_{T} = -\frac{1}{V} \left( \frac{\partial V}{\partial P} \right)_{T} = \frac{V}{N k T} \frac{\Phi_{s-1}(z)}{\Phi_{s}(z)}.
\end{equation*}

沿着绝热线
\begin{equation*}
	&S, \; z, \; P T^{-s-1}, \; V T^{s}, \; P V^{\frac{s+1}{s}} = \text{const}.\\
	&\kappa_{S} = -\frac{1}{V} \left( \frac{\partial V}{\partial P} \right)_{S} = \frac{s}{s+1} \frac{V}{N k T} \frac{\Phi_{s}(z)}{\Phi_{s+1}(z)}. \quad
	c_{s}^{2} = \left( \frac{\partial P}{\partial \rho} \right)_{S} = \frac{V}{m \kappa_{S}} = \frac{s+1}{s} \frac{k T}{m} \frac{\Phi_{s+1}(z)}{\Phi_{s}(z)}.
\end{equation*}

不难给出
\begin{equation*}
	\gamma = \frac{C_{P}}{C_{V}} = \frac{\kappa_{T}}{\kappa_{S}} = \frac{s+1}{s} \frac{\Phi_{s+1}(z) \Phi_{s-1}(z)}{\Phi_{s}^{2}(z)}.
\end{equation*}

\subsubsection{玻色凝聚与费米低温展开}

【对于玻色子】

由于 \(s > 1\) 时 \(g_{s}(z) \le \zeta(s)\)，等号在 \(z = 1\) 取得，因此如上求出的激发态所能容纳的粒子数有限，当温度低于临界值 \(T_{c}\) 时，部分粒子将凝聚到基态，使基态有宏观多的粒子数 \(N_{0}\)，其中
\begin{equation*}
	T_{c} = \frac{1}{k} \left( \frac{N}{B V \Gamma(s+1) \zeta(s)} \right)^{1/s}, \quad N_{0} = N \left( 1 - \left( \frac{T}{T_{c}} \right)^{s} \right).
\end{equation*}

当 \(T < T_{c}\) 时，降温不再改变化学势 （始终为零，对应 \(z = 1\)），而仅仅改变激发态的粒子数量
\begin{equation*}
	S &= (s+1) N k \left( \frac{T}{T_{c}} \right)^{s} \frac{\zeta(s+1)}{\zeta(s)}. \quad
	P = \frac{N k T}{V} \frac{\zeta(s+1)}{\zeta(s)} \left( \frac{T}{T_{c}} \right)^{s}. \\
	C_{V} &= T \left( \frac{\partial S}{\partial T} \right)_{V} = s (s+1) N k \left( \frac{T}{T_{c}} \right)^{s} \frac{\zeta(s+1)}{\zeta(s)}. \quad
	\left( \frac{\partial C_{V}}{\partial T} \right)_{V} = s^{2} (s+1) \frac{N k}{T_{c}} \left( \frac{T}{T_{c}} \right)^{s-1} \frac{\zeta(s+1)}{\zeta(s)}.
\end{equation*}

在相变点 \(T = T_{c}\) 附近
\begin{equation*}
	&S = (s+1) N k \frac{\zeta(s+1)}{\zeta(s)},\quad C_{V}\big|_{T_{c}^{-}} = s (s+1) N k \frac{\zeta(s+1)}{\zeta(s)},\quad \left( \frac{\partial C_{V}}{\partial T} \right)_{V} = s^{2} (s+1) \frac{\zeta(s+1)}{\zeta(s)} \frac{N k}{T_{c}}.\\
	&\Delta C_{V} = C_{V} \big|_{T_{c}^{-}}^{T_{c}^{+}} =
	\begin{cases}
		-s^{2} N k \dfrac{\zeta(s)}{\zeta(s-1)}, & s > 2.   \\
		0,                                       & s \le 2.
	\end{cases}\\
	&\Delta \left.\left( \frac{\partial C_{V}}{\partial T} \right)\right|_{V} = \left( \frac{\partial C_{V}}{\partial T} \right)_{V} \big|_{T_{c}^{-}}^{T_{c}^{+}} =
	\begin{cases}
		\displaystyle -s^{2} \frac{N k}{T_{c}} \left( \frac{\zeta(s)}{\zeta(s-1)} + s \frac{\zeta^{2}(s) \zeta(s-2)}{\zeta^{3}(s-1)} \right), & s \ge \dfrac{3}{2}. \\
		0,                                                                                                                                    & s < \dfrac{3}{2}.
	\end{cases}
\end{equation*}

可见，\(S\)总是连续的，\(C_{V}\) 通常是突变的 （除非 \(s \le 2\)），\(\left(\partial C_{V}/\partial T\right)_{V}\) 也通常是突变的 （除非 \(s < 3/2\)）.

【对于费米子】

温度很低时，粒子主要分布在\(\varepsilon \le \varepsilon_{F}\) 的费米球内，其中
\begin{equation*}
	\varepsilon_{F} = \left( \frac{N}{B V} \right)^{1/s}.
\end{equation*}

各热学量为
\begin{equation*}
	\mu &\simeq \varepsilon_{F} \left( 1 - \frac{s-1}{6} \left( \frac{\pi k T}{\varepsilon_{F}} \right)^{2} - \frac{s-1}{360} (57 - 35 s + 2 s^{2}) \left( \frac{\pi k T}{\varepsilon_{F}} \right)^{4} \right). \\
	P &\simeq \frac{1}{s+1} \frac{N}{V} \varepsilon_{F} \left( 1 + \frac{1}{6} (s+1) \left( \frac{\pi k T}{\varepsilon_{F}} \right)^{2} + \frac{(s+1)(57 - 53 s + 6 s^{2})}{360} \left( \frac{\pi k T}{\varepsilon_{F}} \right)^{4} \right). \\
	E &\simeq \frac{s}{s+1} N \varepsilon_{F} \left( 1 + \frac{1}{6} (s+1) \left( \frac{\pi k T}{\varepsilon_{F}} \right)^{2} + \frac{(s+1)(57 - 53 s + 6 s^{2})}{360} \left( \frac{\pi k T}{\varepsilon_{F}} \right)^{4} \right). \\
	C_{V} &\simeq S \simeq \frac{s \pi^{2}}{3} N k \frac{k T}{\varepsilon_{F}} \left( 1 + \frac{57 - 53 s + 6 s^{2}}{30} \pi^{2} \left( \frac{k T}{\varepsilon_{F}} \right)^{2} \right).
\end{equation*}

\begin{example}
	简并度为\(g\)的\(d\)维非相对论粒子
	\begin{equation*}
		\varepsilon = \frac{p^{2}}{2m}, \quad \Sigma(\varepsilon) = g \frac{2 \pi^{d/2} V}{d\; \Gamma(d/2) h^{d}} (2m \varepsilon)^{d/2}, \quad B = g \left( \frac{2 \pi m}{h^{2}} \right)^{d/2} \frac{1}{\Gamma(d/2+1)}, \quad s = \frac{d}{2}.
	\end{equation*}

	若是玻色子，给出凝聚温度
	\begin{equation*}
		T_c=\frac{h^2}{2\pi m k}\left(\frac{N}{gV\zeta(d/2)}\right)^{2/d}.
	\end{equation*}

	若是给出费米能级、费米波矢、零温简并压为
	\begin{equation*}
    \varepsilon_{F} &= \frac{h^{2}}{2\pi m} \left( \frac{N \Gamma(d/2+1)}{V g} \right)^{2/d},\quad
    k_{F} = \left( \frac{2^{d} \Gamma(d/2+1) \pi^{d/2}}{g} \frac{N}{V} \right)^{1/d},\\
    P_0 &= \frac{h^{2}}{m} \frac{1}{\pi (d+2)} \left( \frac{\Gamma(d/2+1)}{g} \right)^{2/d} \left( \frac{N}{V} \right)^{2/d + 1}.
	\end{equation*}
\end{example}

\subsection{占据数统计：弱相互作用计算的统计预备}

理想气体的热力学量已经给出基准解，但相互作用修正还需要更高阶的占据数矩和涨落。把这一部分单独列出，可以避免在后面的二阶微扰计算中反复插入统计恒等式，也能明确区分“理想气体统计”与“相互作用矩阵元”.

对于任意关于状态 \(\alpha\) (不含自旋，自旋简并度单独记为\(g\)) 的物理量
\begin{equation*}
	\langle X_{\alpha} \rangle =
	\begin{cases}
		\displaystyle \frac{ \sum_{n=0}^{\infty} X[n] z^{n} e^{-n \beta \varepsilon_{\alpha}}}{ \sum_{n=0}^{\infty} z^{n} e^{-n \beta \varepsilon_{\alpha}}} = \sum_{n=0}^{\infty} X[n] z^{n} e^{-n \beta \varepsilon_{\alpha}} \bigl( 1 - z e^{-\beta \varepsilon_{\alpha}} \bigr), & \text{B.E.} \\[12pt]
		\displaystyle \frac{X[0] + z e^{-\beta \varepsilon_{\alpha}} X[1]}{1 + z e^{-\beta \varepsilon_{\alpha}}},                                                                                                                                                                   & \text{F.D.}
	\end{cases}
\end{equation*}

比如
\begin{equation*}
	&\langle n_{\alpha} \rangle =
	\begin{cases}
		\dfrac{1}{z^{-1} e^{\beta \varepsilon_{\alpha}} - 1}, & \text{B.E.} \\[8pt]
		\dfrac{1}{z^{-1} e^{\beta \varepsilon_{\alpha}} + 1}, & \text{F.D.}
	\end{cases}\quad\mathrm{Var}(n_{\alpha}) = \langle n_{\alpha}^{2} \rangle - \langle n_{\alpha} \rangle^{2} =
	\begin{cases}
		\langle n_{\alpha} \rangle (1 + \langle n_{\alpha} \rangle), & \text{B.E.} \\
		\langle n_{\alpha} \rangle (1 - \langle n_{\alpha} \rangle), & \text{F.D.}
	\end{cases}\\
	&\mathrm{Var}(n_{\alpha}^{2}) = \langle n_{\alpha}^{4} \rangle - \langle n_{\alpha}^{2} \rangle^{2} =
	\begin{cases}
		\langle n_{\alpha} \rangle \left( 1 + 13 \langle n_{\alpha} \rangle + 32 \langle n_{\alpha} \rangle^{2} + 20 \langle n_{\alpha} \rangle^{3} \right), & \text{B.E.} \\
		\langle n_{\alpha} \rangle (1 - \langle n_{\alpha} \rangle),                                                                                         & \text{F.D.}
	\end{cases}\\
	&\mathrm{Cov}(n_{\alpha}, n_{\alpha}^{2}) = \langle n_{\alpha}^{3} \rangle - \langle n_{\alpha} \rangle \langle n_{\alpha}^{2} \rangle =
	\begin{cases}
		\langle n_{\alpha} \rangle (1 + \langle n_{\alpha} \rangle) (1 + 4 \langle n_{\alpha} \rangle), & \text{B.E.} \\
		\langle n_{\alpha} \rangle (1 - \langle n_{\alpha} \rangle),                                    & \text{F.D.}
	\end{cases}
\end{equation*}

为避免重复计入自旋简并度，本小节先按单一自旋分量的动量态求和，记
\begin{equation*}
	\mathcal{N} = \sum_{\alpha} \langle n_{\alpha} \rangle, \quad \delta \mathcal{N} = \sum_{\alpha} \left( n_{\alpha} - \langle n_{\alpha} \rangle \right).
\end{equation*}

若单自旋态密度写作 \(\Sigma_0(\varepsilon)=B_0V\varepsilon^s\)，而上一节的总态密度系数为 \(B=gB_0\)，则求和连续化为
\begin{equation*}
	&\mathcal{N} = \sum_{\alpha} \frac{1}{z^{-1} e^{\beta \varepsilon_{\alpha}} \mp 1} = \int_{0}^{\infty} \frac{s B_0 V \varepsilon^{s-1}}{z^{-1} e^{\beta \varepsilon} \mp 1} \, \d \varepsilon = B_0 V \Gamma(s+1) (kT)^{s} \Phi_{s}(z).\\
	&\langle (\delta \mathcal{N})^{2} \rangle = B_0 V \Gamma(s+1) (kT)^{s} \Phi_{s-1}(z).\quad
	\langle (\delta \mathcal{N})^{3} \rangle = B_0 V \Gamma(s+1) (kT)^{s} \Phi_{s-2}(z).\\
	&\langle (\delta \mathcal{N})^{4} \rangle - 3 \langle (\delta \mathcal{N})^{2} \rangle^{2} = B_0 V \Gamma(s+1) (kT)^{s} \Phi_{s-3}(z).
\end{equation*}
总粒子数为 \(N=g\mathcal{N}\)，等价地可继续使用已经吸收简并度的 \(B\)；后面的弱相互作用矩阵元把自旋组合因子显式抽出，所以动量求和中不再额外乘 \(g\).

根据统计数学的定义，不难给出
\begin{equation*}
	\left\langle 2 \sum_{\alpha \neq \beta} n_{\alpha} n_{\beta} + \sum_{\alpha} n_{\alpha} (n_{\alpha} - 1) \right\rangle
	&= \langle \mathcal{N} \rangle^{2} - \langle \mathcal{N} \rangle + \left\langle (\delta \mathcal{N})^{2} \right\rangle. \\
	\mathrm{Var}\!\left( 2 \sum_{\alpha \neq \beta} n_{\alpha} n_{\beta} + \sum_{\alpha} n_{\alpha} (n_{\alpha} - 1) \right)
	&= 4 (\langle \mathcal{N} \rangle^2 - \langle \mathcal{N} \rangle) \left\langle (\delta \mathcal{N})^{2} \right\rangle
	+ \left\langle (\delta \mathcal{N})^{4} \right\rangle+(4 \langle \mathcal{N} \rangle - 2) \left\langle (\delta \mathcal{N})^{3} \right\rangle
	\\&\qquad- \left\langle (\delta \mathcal{N})^{2} \right\rangle \left( \left\langle (\delta \mathcal{N})^{2} \right\rangle - 1 \right).
\end{equation*}

在热力学极限下，保留到体积的最大幂次，得到
\begin{equation*}
	&\left\langle 2 \sum_{\alpha \neq \beta} n_{\alpha} n_{\beta} + \sum_{\alpha} n_{\alpha} (n_{\alpha} - 1) \right\rangle \simeq \left( \sum_{\alpha} \langle n_{\alpha} \rangle \right)^{2}.\\
	&\mathrm{Var}\left( 2 \sum_{\alpha \neq \beta} n_{\alpha} n_{\beta} + \sum_{\alpha} n_{\alpha} (n_{\alpha} - 1) \right) \simeq 4 \left( \sum_{\alpha} \langle n_{\alpha} \rangle \right)^{2} \sum_{\alpha} \left( \langle n_{\alpha}^{2} \rangle - \langle n_{\alpha} \rangle^{2} \right).
\end{equation*}

特别地，对于费米子
\begin{equation*}
	2 \sum_{\alpha \neq \beta} n_{\alpha} n_{\beta} + \sum_{\alpha} n_{\alpha} (n_{\alpha} - 1) = 2 \sum_{\alpha \neq \beta} n_{\alpha} n_{\beta}.
\end{equation*}

\subsection{弱相互作用修正：巨配分函数的二阶展开}

在本节中，\(\mathcal{H}_{1}\) 由前面的低能接触赝势给出。修正按 \(a/\lambda\) 的幂次组织。下述展开只在非凝聚态和热力学极限的指定阶次上使用；玻色凝聚态不应直接代入这一微扰级数.

\subsubsection{正则配分函数与巨配分函数的形式展开}

\(N\) 粒子配分函数为
\begin{equation*}
	Q_{N} = \mathrm{Tr} \, e^{-\beta \mathcal{H}} = \mathrm{Tr} \, e^{-\beta (\mathcal{H}_{0} + \mathcal{H}_{1})}.
\end{equation*}

在非凝聚态下，配分函数可以展开到硬球势参量的零阶、一阶、二阶 (小量为 \(a/\lambda\) 量级，\(\lambda\) 是热波长)
\begin{equation*}
	Q_{N}^{(0)} &= \sum_{r} e^{-\beta E_{r}}, \\
	Q_{N}^{(1)} &= \sum_{k=1}^{\infty} \frac{(-\beta)^{k}}{k!} \mathrm{Tr} \sum_{l=0}^{k-1} H_{0}^{l} H_{1} H_{0}^{k-l-1} = -\beta \mathrm{Tr} \, e^{-\beta H_{0}} H_{1} = -\beta \sum_{r} e^{-\beta E_{r}} H_{rr}. \\
	Q_{N}^{(2)} &= \sum_{k=2}^{\infty} \frac{(-\beta)^{k}}{k!} \mathrm{Tr} \sum_{\substack{p,q,s=0 \\ p+q+s=k-2}} H_{0}^{p} H_{1} H_{0}^{q} H_{1} H_{0}^{s} = \beta^{2} \sum_{i=0}^{\infty} \sum_{j=0}^{\infty} \frac{(j+1)(-\beta)^{i+j}}{(i+j+2)!} \sum_{r,s} E_{r}^{j} E_{s}^{i} |H_{rs}|^{2} \\
	&= \frac{1}{2} \beta^{2} \sum_{r} e^{-\beta E_{r}} |H_{rr}|^{2} + \beta \sum_{r \neq s} \frac{|H_{rs}|^{2}}{E_{r} - E_{s}} e^{-\beta E_{s}}.
\end{equation*}
式中，\(r, s\)代表 \(\{n_{\alpha}\}\) 分布的序号数，\(H_{rs}\) 是 \(H_{1}\) 在相应分布上的矩阵元，粒子数满足约束 \(N = \sum_{\alpha} n_{\alpha}\)，能量可表示为 \(E = \sum_{\alpha} n_{\alpha} \varepsilon_{\alpha}\)，计算中用到
\begin{equation*}
	\sum_{i=0}^{\infty} \sum_{j=0}^{\infty} \frac{(j+1) x^{j} y^{i}}{(i+j+2)!} = \frac{e^{y} - e^{x} - (y-x)e^{x}}{(y-x)^{2}}.
\end{equation*}

同样地，可以将巨配分函数展开
\begin{equation*}
	\Xi &= \sum_{N=0}^{\infty} z^{N} Q_{N}, \quad z = e^{\beta\mu}.\\
	\Xi^{(0)} &= \sum_{N=0}^{\infty} z^{N} \sum_{r} e^{-\beta E_{r}} = \prod_{\alpha} \sum_{n_{\alpha}} \left( z e^{-\beta \varepsilon_{\alpha}} \right)^{n_{\alpha}}. \\
	\Xi^{(1)} &= -\beta \Xi^{(0)} \langle H_{1} \rangle. \quad
	\Xi^{(2)} = \frac{1}{2} \beta^{2} \Xi^{(0)} \langle |H_{1}|^{2} \rangle + \beta \Xi^{(0)} \langle H_{1} \mathcal{R} H_{1} \rangle.
\end{equation*}
其中，\(\langle \cdot \rangle\) 指以 \(z e^{-\beta E_{r}}\) 为权重 (须归一化) 取系综平均，\(\mathcal{R}\)作用到算符\(A\)上将得到
\begin{equation*}
	(\mathcal{R} A)_{rs} = \frac{A_{rs}}{E_{r} - E_{s}} (1 - \delta_{rs}).
\end{equation*}

进而，常用的巨配分函数对数可以展开为
\begin{equation*}
	\ln \Xi &= \ln \Xi^{(0)} + \frac{\Xi^{(1)}}{\Xi^{(0)}} + \left( \frac{\Xi^{(2)}}{\Xi^{(0)}} - \frac{1}{2} \left( \frac{\Xi^{(1)}}{\Xi^{(0)}} \right)^{2} \right) \\
	&= \sum_{\alpha} \ln \left( \sum_{n_{\alpha}} \left( z e^{-\beta \varepsilon_{\alpha}} \right)^{n_{\alpha}} \right) - \beta \langle H_{1} \rangle + \frac{1}{2} \beta^{2} \mathrm{Var}(H_{1}) + \beta \langle H_{1} \mathcal{R} H_{1} \rangle.
\end{equation*}
式中 \(\mathrm{Var}\) 指取系综方差.

\subsubsection{二次量子化、自旋简并与动量守恒}

在此列出玻色子、费米子的产生湮灭算符满足的关系式，以备后文适用
\begin{equation*}
	&[a_{\alpha}, a_{\beta}^{\dagger}] = \delta_{\alpha\beta}, \quad [a_{\alpha}, a_{\beta}] = [a_{\alpha}^{\dagger}, a_{\beta}^{\dagger}] = 0, \quad \text{B.E.} \\
	&\{a_{\alpha}, a_{\beta}^{\dagger}\} = \delta_{\alpha\beta}, \quad \{a_{\alpha}, a_{\beta}\} = \{a_{\alpha}^{\dagger}, a_{\beta}^{\dagger}\} = 0, \quad \text{F.D.} \\
	&a_{\alpha} \ket{\{n_{1} \ldots n_{\alpha} \ldots\}} = \sqrt{n_{\alpha}} \ket{\{n_{1} \ldots n_{\alpha} - 1 \ldots\}}, \\
	&a_{\alpha}^{\dagger} \ket{\{n_{1} \ldots n_{\alpha} \ldots\}} = \sqrt{n_{\alpha} + 1} \ket{\{n_{1} \ldots n_{\alpha} + 1 \ldots\}}, \quad
	n_{\alpha} = a_{\alpha}^{\dagger} a_{\alpha}.
\end{equation*}

将硬球势量子化得
\begin{equation*}
	H_{1} = \frac{2\pi a \hbar^{2}}{m V} \sum_{\alpha,\beta,\gamma,\lambda} a_{\alpha}^{\dagger} a_{\beta}^{\dagger} a_{\gamma} a_{\lambda} \delta_{\vb{p}_{\alpha} + \vb{p}_{\beta}, \vb{p}_{\gamma} + \vb{p}_{\lambda}}.
\end{equation*}
其中 \(\alpha\) 包含动量 \(\vb{p}\) 和自旋 \(\sigma\) 两个状态，\(\delta\,\)(Kronecker)符号给出动量守恒的约束.

先考虑对自旋的求和.

由于赝势对应的s波 (\(l = 0\)) 空间波函数是对称的，因此玻色子的自旋波函数须对称，费米子的自旋波函数须反对称。又考虑到协方差的非零项必然要求某个自旋共用，故对于自旋为\(J\) (简并度\(g = 2J + 1\)) 的粒子而言，巨配分函数各项的自旋简并因子为
\begin{equation*}
	\langle H_{1} \rangle \propto
	\begin{cases}
		g(g+1), & \text{B.E.} \\
		g(g-1), & \text{F.D.}
	\end{cases}
	\quad \mathrm{Var}(H_{1}) \propto
	\begin{cases}
		g(g+1)^{2}, & \text{B.E.} \\
		g(g-1)^{2}, & \text{F.D.}
	\end{cases}\quad
	\langle H_{1} \mathcal{R} H_{1} \rangle \big|_{(1,2) \to (3,4)} \propto
	\begin{cases}
		2g(g+1), & \text{B.E.} \\
		2g(g-1), & \text{F.D.}
	\end{cases}
\end{equation*}
其中下标 \((1,2) \to (3,4)\) 代表由两种不同动量态散射至另外两种不同动量态的情形.

再考虑动量部分，以下 \(\alpha, \beta\) 只代表动量 \(\vb{p}\) 态.

\subsubsection{玻色子贡献}

在粒子数表象下，哈密顿量的对角元给出
\begin{equation*}
	\sum_{\alpha,\beta,\gamma,\lambda} a_{\alpha}^{\dagger} a_{\beta}^{\dagger} a_{\gamma} a_{\lambda} \Big|_{\text{diag}}= 2 \sum_{\alpha < \beta} n_{\alpha} n_{\beta} + \sum_{\alpha} n_{\alpha} (n_{\alpha} - 1).
\end{equation*}

取系综平均，在热力学极限下，保留关于体积的最大贡献项得
\begin{equation*}
	&\left\langle 2 \sum_{\alpha < \beta} n_{\alpha} n_{\beta} + \sum_{\alpha} n_{\alpha} (n_{\alpha} - 1) \right\rangle \simeq \left( \sum_{\alpha} \langle n_{\alpha} \rangle \right)^{2} = \left( \frac{V}{\lambda^{3}} g_{3/2}(z) \right)^{2}.\\
	&\mathrm{Var}\left( 2 \sum_{\alpha < \beta} n_{\alpha} n_{\beta} + \sum_{\alpha} n_{\alpha} (n_{\alpha} - 1) \right) \simeq 4 \left( \sum_{\alpha} \langle n_{\alpha} \rangle \right)^{2} \sum_{\alpha} \left( \langle n_{\alpha}^{2} \rangle - \langle n_{\alpha} \rangle^{2} \right) = \frac{4V^{3}}{\lambda^{9}} (g_{3/2}(z))^{2} g_{1/2}(z).
\end{equation*}

因此
\begin{equation*}
	\langle H_{1} \rangle &= \left( \frac{V}{\lambda^{3}} g_{3/2}(z) \right)^{2} \cdot \frac{2\pi a \hbar^{2}}{m V} \cdot 2(J+1)(2J+1) = \frac{(2J+1) V}{\beta \lambda^{3}} \frac{2(J+1) a}{\lambda} (g_{3/2}(z))^{2}. \\
	\mathrm{Var}(H_{1}) &= \frac{4V^{3}}{\lambda^{9}} (g_{3/2}(z))^{2} g_{1/2}(z) \cdot \left( \frac{2\pi a \hbar^{2}}{m V} \right)^{2} \cdot 4(J+1)^{2}(2J+1) = \frac{(2J+1) V}{\beta^{2} \lambda^{3}} \frac{16(J+1)^{2} a^{2}}{\lambda^{2}} g_{1/2}(z) (g_{3/2}(z))^{2}.
\end{equation*}

而 \(\langle H_{1} \mathcal{R} H_{1} \rangle\) 项对应着 \((1,2) \to (3,4)\), \((1,1) \to (2,3)\), \((1,2) \to (3,3)\) 三类二体碰撞情况，每种情况均满足动量守恒 \(\vb{p}_{\alpha} + \vb{p}_{\beta} = \vb{p}_{\gamma} + \vb{p}_{\lambda}\)。由于后两种情况固定死了一个粒子，因此相较第一项，求和时会少一个体积因子，在热力学极限下，贡献可忽略。故仅保留第一种二体碰撞，得到
\begin{equation*}
	\langle H_{1} \mathcal{R} H_{1} \rangle &= 4(J+1)(2J+1) \cdot \left( \frac{2\pi a \hbar^{2}}{m V} \right)^{2} \sum_{\vb{p}_{\alpha},\vb{p}_{\beta},\vb{p}_{\gamma},\vb{p}_{\lambda}} \frac{\langle n_{\gamma} \rangle \langle n_{\lambda} \rangle (1 + \langle n_{\alpha} \rangle) (1 + \langle n_{\beta} \rangle)}{\varepsilon_{\alpha} + \varepsilon_{\beta} - \varepsilon_{\gamma} - \varepsilon_{\lambda}} \delta_{\vb{p}_{\alpha} + \vb{p}_{\beta}, \vb{p}_{\gamma} + \vb{p}_{\lambda}} \\
	&= 8(J+1)(2J+1) \left( \frac{2\pi a \hbar^{2}}{m V} \right)^{2} \sum_{\vb{p}_{\alpha},\vb{p}_{\beta},\vb{p}_{\gamma},\vb{p}_{\lambda}} \frac{\langle n_{\gamma} \rangle \langle n_{\lambda} \rangle \langle n_{\alpha} \rangle}{\varepsilon_{\alpha} + \varepsilon_{\beta} - \varepsilon_{\gamma} - \varepsilon_{\lambda}} \delta_{\vb{p}_{\alpha} + \vb{p}_{\beta}, \vb{p}_{\gamma} + \vb{p}_{\lambda}}.
\end{equation*}
其中由于求和的对称性，四个及两个 \(\langle n_{\alpha} \rangle\) 项的求和为零.

进一步换元，并将求和化为积分
\begin{equation*}
	\vb{p}_{\alpha,\beta} = \frac{\vb{\kappa}}{2} \pm \vb{q}, \quad \vb{p}_{\gamma,\lambda} = \frac{\vb{\kappa}}{2} \pm \vb{p}, \quad \sum_{\alpha} \to \frac{V}{h^{3}} \int \d \vb{p}_{\alpha}, \quad \d \vb{p}_{\alpha} \d \vb{p}_{\beta} \d \vb{p}_{\gamma} = \d \vb{\kappa} \, \d \vb{p} \, \d \vb{q}.
\end{equation*}
代入
\begin{equation*}
	\langle n_{\alpha} \rangle = \frac{1}{z^{-1} e^{\beta \varepsilon_{\alpha}} - 1} = \sum_{l=1}^{\infty} z^{l} e^{-l\beta \varepsilon_{\alpha}}, \quad \varepsilon_{\alpha} + \varepsilon_{\beta} - \varepsilon_{\gamma} - \varepsilon_{\lambda} = \frac{q^{2} - p^{2}}{m}.
\end{equation*}
积分得
\begin{equation*}
	\langle H_{1} \mathcal{R} H_{1} \rangle &= 8(J+1)(2J+1) \left( \frac{2\pi a \hbar^{2}}{m V} \right)^{2} m \left( \frac{V}{h^{3}} \right)^{3} \sum_{j=1}^{\infty} \sum_{k=1}^{\infty} \sum_{l=1}^{\infty} \iiint \frac{z^{j+k+l}}{q^{2} - p^{2}}\\&\qquad \exp\left[ -\frac{\lambda^{2}}{4\pi\hbar^{2}} \left( \frac{j+k+l}{4} \kappa^{2} + \vb{\kappa} \cdot (j\vb{q} + (k-l)\vb{p}) + (j q^{2} + (k+l) p^{2}) \right) \right] \d \vb{\kappa} \, \d \vb{p} \, \d \vb{q} \\
	&= \frac{(2J+1) V}{\beta \lambda^{3}} \frac{8(J+1) a^{2}}{\lambda^{2}} G(z).
\end{equation*}
式中
\begin{equation*}
	G(z) = \sum_{j=1}^{\infty} \sum_{k=1}^{\infty} \sum_{l=1}^{\infty} \frac{z^{j+k+l}}{(k+l)(j+l) \sqrt{jkl}}.
\end{equation*}
这个三重级数在玻色气体的物理区间 \(0\le z\le1\) 绝对收敛；接近凝聚点时只能从 \(z<1\) 取极限，不能把它和凝聚态本身的微扰展开混同.
并用到公式
\begin{equation*}
	\int_{0}^{\infty} \int_{0}^{\infty} \frac{\sinh(u p q)}{q^{2} - p^{2}} p q e^{-v q^{2} - w p^{2}} \, \d p \, \d q = \frac{\pi}{4} \frac{u (w - v)}{((w+v)^{2} - u^{2}) \sqrt{4wv - u^{2}}}.
\end{equation*}

综合以上各项知
\begin{equation*}
	\ln \Xi = \frac{(2J+1) V}{\lambda^{3}} \left( g_{5/2}(z) - \frac{2(J+1) a}{\lambda} (g_{3/2}(z))^{2} + \frac{8(J+1)^{2} a^{2}}{\lambda^{2}} g_{1/2}(z) (g_{3/2}(z))^{2} + \frac{8(J+1) a^{2}}{\lambda^{2}} G(z) \right).
\end{equation*}

\subsubsection{费米子贡献}

由于同一个态只能占据一个粒子 (\(n_{\alpha} = 0, 1\))，因此在粒子数表象下，哈密顿量的对角元给出
\begin{equation*}
	\sum_{\alpha,\beta,\gamma,\lambda} a_{\alpha}^{\dagger} a_{\beta}^{\dagger} a_{\gamma} a_{\lambda}\Big|_{\text{diag}} = 2 \sum_{\alpha < \beta} n_{\alpha} n_{\beta}.
\end{equation*}

重复玻色子的操作，仅保留体积的最高次项得
\begin{equation*}
	&\left\langle 2 \sum_{\alpha < \beta} n_{\alpha} n_{\beta} \right\rangle \simeq \left( \sum_{\alpha} \langle n_{\alpha} \rangle \right)^{2} = \left( \frac{V}{\lambda^{3}} f_{3/2}(z) \right)^{2}.\\
	&\mathrm{Var}\left( 2 \sum_{\alpha < \beta} n_{\alpha} n_{\beta} \right) \simeq 4 \left( \sum_{\alpha} \langle n_{\alpha} \rangle \right)^{2} \sum_{\alpha} \left( \langle n_{\alpha}^{2} \rangle - \langle n_{\alpha} \rangle^{2} \right) = \frac{4V^{3}}{\lambda^{9}} (f_{3/2}(z))^{2} f_{1/2}(z).
\end{equation*}

因此
\begin{equation*}
	\langle H_{1} \rangle &= \left( \frac{V}{\lambda^{3}} f_{3/2}(z) \right)^{2} \cdot \frac{2\pi a \hbar^{2}}{m V} \cdot 2J(2J+1) = \frac{(2J+1) V}{\beta \lambda^{3}} \frac{2J a}{\lambda} (f_{3/2}(z))^{2}. \\
	\mathrm{Var}(H_{1}) &= \frac{4V^{3}}{\lambda^{9}} (f_{3/2}(z))^{2} f_{1/2}(z) \cdot \left( \frac{2\pi a \hbar^{2}}{m V} \right)^{2} \cdot 4J^{2}(2J+1) = \frac{(2J+1) V}{\beta^{2} \lambda^{3}} \frac{16J^{2} a^{2}}{\lambda^{2}} f_{1/2}(z) (f_{3/2}(z))^{2}.
\end{equation*}

注意到费米子的产生算符等价地满足
\begin{equation*}
	a_{\alpha}^{\dagger} \ket{\{n_{1} \ldots n_{\alpha} \ldots\}} = \sqrt{1 - n_{\alpha}} \ket{\{n_{1} \ldots n_{\alpha} + 1 \ldots\}}.
\end{equation*}
因此只存在 \((1,2) \to (3,4)\) 一种二体碰撞情形，对应
\begin{equation*}
	\langle H_{1} \mathcal{R} H_{1} \rangle &= 4J(2J+1) \cdot \left( \frac{2\pi a \hbar^{2}}{m V} \right)^{2} \sum_{\vb{p}_{\alpha},\vb{p}_{\beta},\vb{p}_{\gamma},\vb{p}_{\lambda}} \frac{\langle n_{\gamma} \rangle \langle n_{\lambda} \rangle (1 - \langle n_{\alpha} \rangle) (1 - \langle n_{\beta} \rangle)}{\varepsilon_{\alpha} + \varepsilon_{\beta} - \varepsilon_{\gamma} - \varepsilon_{\lambda}} \delta_{\vb{p}_{\alpha} + \vb{p}_{\beta}, \vb{p}_{\gamma} + \vb{p}_{\lambda}} \\
	&= -8 J (2J+1) \left( \frac{2\pi a \hbar^{2}}{m V} \right)^{2} \sum_{\vb{p}_{\alpha},\vb{p}_{\beta},\vb{p}_{\gamma},\vb{p}_{\lambda}} \frac{\langle n_{\gamma} \rangle \langle n_{\lambda} \rangle \langle n_{\alpha} \rangle}{\varepsilon_{\alpha} + \varepsilon_{\beta} - \varepsilon_{\gamma} - \varepsilon_{\lambda}} \delta_{\vb{p}_{\alpha} + \vb{p}_{\beta}, \vb{p}_{\gamma} + \vb{p}_{\lambda}}.
\end{equation*}
代入
\begin{equation*}
	\langle n_{\alpha} \rangle = \frac{1}{z^{-1} e^{\beta \varepsilon_{\alpha}} + 1} = -\sum_{l=1}^{\infty} (-z)^{l} e^{-l\beta \varepsilon_{\alpha}}.
\end{equation*}
完全重复玻色子的操作，得到
\begin{equation*}
	\langle H_{1} \mathcal{R} H_{1} \rangle = -\frac{(2J+1) V}{\beta \lambda^{3}} \frac{8J a^{2}}{\lambda^{2}} F(z).
\end{equation*}
式中
\begin{equation*}
	F(z) = \sum_{j=1}^{\infty} \sum_{k=1}^{\infty} \sum_{l=1}^{\infty} \frac{(-1)^{j+k+l-1} z^{j+k+l}}{(k+l)(j+l) \sqrt{jkl}} = -G(-z).
\end{equation*}
这里的三重级数只是 \(|z|\le1\) 内的高温展开；费米强简并区 \(z>1\) 应把 \(F(z)=-G(-z)\) 理解为解析延拓后的定义，后续对 \(\ln z\) 的求导也按这一延拓函数进行.

综合以上各项知
\begin{equation*}
	\ln \Xi = \frac{(2J+1) V}{\lambda^{3}} \left( f_{5/2}(z) - \frac{2J a}{\lambda} (f_{3/2}(z))^{2} + \frac{8J^{2} a^{2}}{\lambda^{2}} f_{1/2}(z) (f_{3/2}(z))^{2} - \frac{8J a^{2}}{\lambda^{2}} F(z) \right).
\end{equation*}

\subsubsection{二阶结果、定粒子数展开与低温费米极限}

【精确到二阶的巨配分函数】

以上讨论给出
\begin{equation*}
	\ln \Xi = \frac{(2J+1) V}{\lambda^{3}}
	\begin{cases}
		g_{5/2}(z) - \dfrac{2(J+1) a}{\lambda} (g_{3/2}(z))^{2} + \dfrac{8(J+1)^{2} a^{2}}{\lambda^{2}} g_{1/2}(z) (g_{3/2}(z))^{2} + \dfrac{8(J+1) a^{2}}{\lambda^{2}} G(z), & \text{B.E.} \\[10pt]
		f_{5/2}(z) - \dfrac{2J a}{\lambda} (f_{3/2}(z))^{2} + \dfrac{8J^{2} a^{2}}{\lambda^{2}} f_{1/2}(z) (f_{3/2}(z))^{2} - \dfrac{8J a^{2}}{\lambda^{2}} F(z),             & \text{F.D.}
	\end{cases}
\end{equation*}

根据热力学公式
\begin{equation*}
	N = z \frac{\partial \ln \Xi}{\partial z}, \quad E = -\frac{\partial \ln \Xi}{\partial \beta} = -kT \frac{\partial \ln \Xi}{\partial \ln \lambda}.
\end{equation*}

按照 \(a/\lambda\) 的幂次展开，保留到二阶项得到
\begin{equation*}
	&E = kT \frac{V}{\lambda^{3}} (2J+1)
	\begin{cases}
		\dfrac{3}{2} g_{5/2}(z_{0}) + 2(J+1) \dfrac{a}{\lambda} g_{3/2}(z_{0})^{2} - 4(J+1) \dfrac{a^{2}}{\lambda^{2}} \left( -5 G(z_{0}) + 3 \dfrac{g_{3/2}(z_{0})}{g_{1/2}(z_{0})} \dfrac{\partial G(z_{0})}{\partial \ln z_{0}} \right), & \text{B.E.} \\[10pt]
		\dfrac{3}{2} f_{5/2}(z_{0}) + 2J \dfrac{a}{\lambda} f_{3/2}(z_{0})^{2} + 4J \dfrac{a^{2}}{\lambda^{2}} \left( -5 F(z_{0}) + 3 \dfrac{f_{3/2}(z_{0})}{f_{1/2}(z_{0})} \dfrac{\partial F(z_{0})}{\partial \ln z_{0}} \right),         & \text{F.D.}
	\end{cases}\\
	&\ln z =
	\begin{cases}
		\ln z_{0} + 4(J+1) g_{3/2}(z_{0}) \dfrac{a}{\lambda} - \dfrac{8(J+1)}{g_{1/2}(z_{0})} \dfrac{\partial G(z_{0})}{\partial \ln z_{0}} \dfrac{a^{2}}{\lambda^{2}}, & \text{B.E.} \\[10pt]
		\ln z_{0} + 4J f_{3/2}(z_{0}) \dfrac{a}{\lambda} + \dfrac{8J}{f_{1/2}(z_{0})} \dfrac{\partial F(z_{0})}{\partial \ln z_{0}} \dfrac{a^{2}}{\lambda^{2}},         & \text{F.D.}
	\end{cases}
\end{equation*}
其中 \(z_{0}\) 是不存在硬球势的逸度，定义为
\begin{equation*}
	N = \frac{(2J+1) V}{\lambda^{3}} \Phi_{3/2}(z_{0}).\quad \Phi_{\nu}(z) =
	\begin{cases}
		g_{\nu}(z), & \text{B.E.} \\
		f_{\nu}(z), & \text{F.D.}
	\end{cases}
\end{equation*}

在零温附近 (\(\ln z \gg 1\))，费米能的修正为
\begin{equation*}
	\mu &= \varepsilon_{F} \left( 1 + \frac{8J}{3\pi} k_{F} a + \frac{8(11 - 2\ln 2) J}{15\pi^{2}} (k_{F} a)^{2} \right) - \frac{k^{2} T^{2}}{\varepsilon_{F}} \frac{\pi^{2}}{12} \left( 1 + \frac{16J}{5\pi^{2}} (7\ln 2 - 1) (k_{F} a)^{2} \right). \\
	E &= n \varepsilon_{F} \left( \frac{3}{5} + \frac{4J}{3\pi} k_{F} a + \frac{8(11 - 2\ln 2) J}{35\pi^{2}} (k_{F} a)^{2} \right) + \frac{n k^{2} T^{2}}{\varepsilon_{F}} \frac{\pi^{2}}{4} \left( 1 + \frac{16J}{15\pi^{2}} (7\ln 2 - 1) (k_{F} a)^{2} \right).
\end{equation*}
其中用到
\begin{equation*}
	&k_{F} = \left( \frac{6\pi^{2} n}{2J+1} \right)^{1/3}, \quad \varepsilon_{F} = \frac{\hbar^{2} k_{F}^{2}}{2m}, \quad \lambda = \hbar \sqrt{\frac{2\pi}{m k T}}, \quad \beta = \frac{1}{kT}.\\
	&F(z) \big|_{\ln z \gg 1} = \frac{16}{105} \frac{11 - 2\ln 2}{\pi^{3/2}} (\ln z)^{7/2} \left( 1 - \frac{35}{8} \frac{2\ln 2 - 1}{11 - 2\ln 2} \frac{\pi^{2}}{(\ln z)^{2}} + \mathcal{O}\left((\ln z)^{-4}\right) \right).
\end{equation*}
这里的强简并展开按 Sommerfeld 型的偶次幂 \( (\ln z)^{-2}\) 组织，因此余项写在括号内.

而由于如上配分函数展开方式不适用于凝聚态，因此无法给出玻色气体在凝聚态的性质.

\subsection{外磁场响应与谱边界诱导的非解析性}

外磁场同时引入 Landau 量子化和 Zeeman 分裂，使得单粒子能级在零温下逐个穿越化学势。高温弱场极限恢复解析的磁化率，而零温有限场则产生由占据边界决定的非解析项.

\subsubsection{一般温度与高温弱场响应}

在外磁场作用下，电荷量为\(q\)、质量为\(m\)的有自旋\(J\)的粒子满足薛定谔方程
\begin{equation*}
	\left(\frac{(\vb{p}-q\vb{A})^{2}}{2m}-\vb{\mu}\cdot\vb{B}\right)\psi
	=i\hbar\frac{\partial\psi}{\partial t}
	=\varepsilon\psi .
\end{equation*}
在\(d\)维空间中，假设外场只作用在\((x_{1},x_{2})\)平面，大小为\(B\,(B>0)\)，磁矢势取为\(\vb{A}=-Bx_{2}\uv{x}_{1}\). 其余\(d-2\)个方向仍是自由的，相应动量记为\(\vb{p}_{\parallel}\). 粒子磁矩满足
\begin{equation*}
	\vb{\mu}=g_{J}\frac{q}{2m}\vb{J},
	\qquad
	J_{z}=m_{J}\hbar,
	\quad
	m_{J}=-J,-J+1,\ldots,J .
\end{equation*}
于是横向运动化为一维谐振子，而沿\(\vb{p}_{\parallel}\)方向保持连续. 能谱为
\begin{equation*}
	\varepsilon_{j,m_{J},\vb{p}_{\parallel}}
	=
	\frac{p_{\parallel}^{2}}{2m}
	+
	\mu_{q}B(2j+1-g_{J}m_{J}),
	\qquad
	\mu_{q}=\frac{|q|\hbar}{2m},
	\quad
	j=0,1,2,\ldots .
\end{equation*}

记\(\sigma=1\)表示玻色子，\(\sigma=-1\)表示费米子，则系统的巨势密度为
\begin{equation*}
	\omega
	&=
	\frac{\sigma}{\beta}
	\frac{|q|B}{h^{d-1}}
	\sum_{j,m_{J}}
	\int_{\mathbb{R}^{d-2}}
	\ln\left[
	1-\sigma z
	\exp\left(
	-\beta\left(
	\frac{p_{\parallel}^{2}}{2m}
	+\mu_{q}B(2j+1-g_{J}m_{J})
	\right)
	\right)
	\right]\d\vb{p}_{\parallel}\\
	&=
	-\frac{1}{\beta\lambda^{d}}
	\sum_{s=1}^{\infty}
	\frac{z^{s}\sigma^{s-1}}{s^{1+d/2}}
	\frac{s\beta\mu_{q}B}{\sinh(s\beta\mu_{q}B)}
	\frac{
		\sinh\!\left(s\beta g_{J}\mu_{q}B(J+1/2)\right)
	}{
		\sinh\!\left(s\beta g_{J}\mu_{q}B/2\right)
	}.
\end{equation*}

当\(\beta\mu_q B\to0\)，即高温弱场极限下，巨势密度近似为
\begin{equation*}
	\omega
	=
	-\frac{2J+1}{\beta\lambda^{d}}
	\left[
		\Phi_{1+d/2}(z)
		+
		\frac{1}{6}
		\left(J(J+1)g_{J}^{2}-1\right)
		\beta^{2}\mu_{q}^{2}B^{2}
		\Phi_{d/2-1}(z)
		+\mathcal{O}(B^{4})
	\right].
\end{equation*}
于是磁化强度为
\begin{equation*}
	M
	=
	-\left(\frac{\partial\omega}{\partial B}\right)_{T,z}
	=
	\frac{2J+1}{3\lambda^{d}}
	\left(J(J+1)g_{J}^{2}-1\right)
	\beta\mu_{q}^{2}B\,\Phi_{d/2-1}(z)
	+\mathcal{O}(B^{3}).
\end{equation*}
可见，磁化是线性的，磁化率为
\begin{equation*}
	\chi_{m}
	=
	\frac{n\mu_{q}^{2}}{3kT}
	\left(J(J+1)g_{J}^{2}-1\right)
	\frac{\Phi_{d/2-1}(z)}{\Phi_{d/2}(z)},
	\qquad
	n=\frac{2J+1}{\lambda^{d}}\Phi_{d/2}(z).
\end{equation*}
式中，\(J(J+1)g_{J}^{2}\)来自 Zeeman 顺磁项，而\(-1\)来自 Landau 轨道抗磁项. 例如电子满足\(J=1/2,g_{J}=2\)，于是顺磁项为抗磁项的3倍.

\subsubsection{零温费米气体与离散占据阈值}

当温度很低时，\(\beta\mu_{q}B\ll1\)不再成立，需要单独求解. 特别考虑零温极限下的费米子，重新积分巨势密度得
\begin{equation*}
	\omega
	&=
	-\frac{1}{\beta}\frac{|q|B}{h^{d-1}}
	\sum_{j,m_{J}}
	\int_{\mathbb{R}^{d-2}}
	\ln\left[
		1+\exp\left(\beta\left(
		\mu-\frac{p_{\parallel}^{2}}{2m}
		-\mu_{q}B(2j+1-g_{J}m_{J})\right)
		\right)
		\right]\d\vb{p}_{\parallel}\\
	&=
	-\frac{|q|B}{h^{d-1}}
	\sum_{j,m_{J}}
	\int_{\mathbb{R}^{d-2}}
	\left(
		\mu-\frac{p_{\parallel}^{2}}{2m}
		-\mu_{q}B(2j+1-g_{J}m_{J})
	\right)_{+}
	\d\vb{p}_{\parallel}\\
	&=
	-\frac{2\mu_{q}B}{\Gamma(d/2+1)}
	\left(\frac{2\pi m}{h^{2}}\right)^{d/2}
	\sum_{j,m_{J}}
	\left(
		\mu-\mu_{q}B(2j+1-g_{J}m_{J})
	\right)_{+}^{d/2}.
\end{equation*}
式中\(x_{+}=\max\{x,0\}=x\Theta(x)\). 用\(l\)表示有序对\((j,m_{J})\)，并记
\begin{equation*}
	\gamma=
	\frac{2\mu_{q}}{\Gamma(d/2)}
	\left(\frac{2\pi m}{h^{2}}\right)^{d/2},
	\qquad
	\eta_{l}=\mu-a_{l}B,
	\qquad
	a_{l}=\mu_{q}(2j+1-g_{J}m_{J}),
	\qquad
	\nu=\frac{d}{2}-1 .
\end{equation*}
则
\begin{equation*}
	\omega
	=
	-\frac{\gamma B}{\nu+1}
	\sum_{l}\eta_{l,+}^{\nu+1},
	\qquad
	n
	=
	-\left(\frac{\partial\omega}{\partial\mu}\right)_{B}
	=
	\gamma B\sum_{l}\eta_{l,+}^{\nu}.
\end{equation*}

给定粒子数密度\(n\)，其约束方程将给出\(\mu\)关于\(B\)的函数. 改变外磁场\(B\)，由于\(\eta_{l,+}\)的阶跃性，会在离散临界值处占据或退出某一有序对. 设\(B=B_{c}\)时序号为\(0\)的有序对恰好被占据，称其为\(0\)号态；除\(0\)号态外所有被占据的\(l\)组成集合\(\mathcal{R}\)，即
\begin{equation*}
	\eta_{0}(B_{c})=0,
	\qquad
	\eta_{l}(B_{c})>0,\quad l\in\mathcal{R}.
\end{equation*}
记
\begin{equation*}
	\xi_{l}=(a_{0}-a_{l})B_{c},
	\qquad
	S_{\alpha}=\sum_{l\in\mathcal{R}}\xi_{l}^{\alpha}.
\end{equation*}
则\(0\)号态占据的临界磁场由
\begin{equation*}
	n
	=
	\gamma B_{c}S_{\nu}
	=
	\gamma B_{c}^{\nu+1}
	\sum_{l\in\mathcal{R}}(a_{0}-a_{l})^{\nu}
\end{equation*}
给出，即
\begin{equation*}
	B_{c}
	=
	\left[
		\frac{n}{\gamma\sum_{l\in\mathcal{R}}(a_{0}-a_{l})^{\nu}}
	\right]^{1/(\nu+1)}.
\end{equation*}

定义辅助化学势\(\bar{\mu}(B)\)满足
\begin{equation*}
	n
	=
	\gamma B
	\sum_{l\in\mathcal{R}}
	\left(\bar{\mu}(B)-a_{l}B\right)^{\nu},
\end{equation*}
而实际化学势满足
\begin{equation*}
	n
	=
	\gamma B
	\sum_{l}
	\left(\mu(B)-a_{l}B\right)_{+}^{\nu}.
\end{equation*}
在\(B_{c}\)附近展开，不难给出
\begin{equation*}
	\bar{\mu}(B_{c})=\mu(B_{c}),
	\qquad
	\bar{\mu}'(B_{c})
	=
	\frac{\sum_{l\in\mathcal{R}}a_{l}\xi_{l}^{\nu-1}}{S_{\nu-1}}
	-\frac{S_{\nu}}{\nu B_{c}S_{\nu-1}}.
\end{equation*}
记
\begin{equation*}
	\eta(B)=\mu(B)-a_{0}B,
	\qquad
	t(B)=\bar{\mu}(B)-a_{0}B .
\end{equation*}
这里\(t(B)\)正比于\(B\)在临界点附近的相对变化率：
\begin{equation*}
	t(B)
	&=
	\left(\bar{\mu}'(B_{c})-a_{0}\right)(B-B_{c})
	+\mathcal{O}\!\left((B-B_{c})^{2}\right)\\
	&=
	-\frac{\nu+1}{\nu}
	\frac{S_{\nu}}{S_{\nu-1}}
	\frac{B-B_{c}}{B_{c}}
	+\mathcal{O}\!\left((B-B_{c})^{2}\right).
\end{equation*}
同理，将粒子数在\(B_c\)附近展开给出
\begin{equation*}
	&t
	=
	\eta
	+
	\frac{\eta_{+}^{\nu}}{\nu S_{\nu-1}}
	+
	\mathcal{O}\!\left(
	(t-\eta)^{2},
	(B-B_{c})(\eta-t),
	(B-B_{c})\eta_{+}^{\nu}
	\right).\\
	&\eta
	\simeq
	t(1-\Theta(t))
	+\Theta(t)
	\begin{cases}
		(\nu S_{\nu-1})^{1/\nu}t^{1/\nu},
		&0<\nu<1,\\[4pt]
		\dfrac{S_{0}}{S_{0}+1}t,
		&\nu=1,\\[8pt]
		t-\frac{t^{\nu}}{\nu S_{\nu-1}},
		&\nu>1.
	\end{cases}
\end{equation*}
可见，当\(t<0\)，即\(B>B_{c}\)时，\(\eta<0\)；当\(t>0\)，即\(B<B_{c}\)时，\(\eta>0\). 这说明，当磁场由大到小跨过\(B_{c}\)时，\(0\)号态由未占据变成被占据，系统热力学函数不解析，出现相变.

\subsubsection{自由能的奇异项与磁化率}

具体从自由能密度角度分析. 系统自由能密度为
\begin{equation*}
	f(B)
	=
	\omega(\mu,B)-\mu\frac{\partial\omega(\mu,B)}{\partial\mu}
	=
	\omega_{\mathcal{R}}(\mu,B)
	-\frac{\gamma B}{\nu+1}\eta_{+}^{\nu+1}
	+\mu n,
\end{equation*}
式中\(\omega_{\mathcal{R}}\)是除去\(0\)号态的巨势密度
\begin{equation*}
	\omega_{\mathcal{R}}(\mu,B)
	=
	-\frac{\gamma B}{\nu+1}
	\sum_{l\in\mathcal{R}}
	(\mu-a_{l}B)^{\nu+1}.
\end{equation*}
而除去\(0\)号态的自由能密度为
\begin{equation*}
	\bar{f}(B)
	=
	\omega_{\mathcal{R}}(\bar{\mu},B)
	-\bar{\mu}
	\frac{\partial\omega_{\mathcal{R}}(\bar{\mu},B)}{\partial\bar{\mu}}
	=
	\omega_{\mathcal{R}}(\bar{\mu},B)+\bar{\mu}n .
\end{equation*}
将\(\omega_{\mathcal{R}}\)在\(\bar{\mu}\)处展开，并利用\(\mu-\bar{\mu}=\eta-t\)，得到相变点附近的奇异项
\begin{equation*}
	f_{\mathrm{sing}}
	&=
	f(B)-\bar{f}(B)\simeq
	\frac{1}{2}
	\left.
	\frac{\partial^{2}\omega_{\mathcal{R}}}{\partial\bar{\mu}^{2}}
	\right|_{B=B_{c}}
	(\eta-t)^{2}
	-\frac{\gamma B_{c}}{\nu+1}\eta_{+}^{\nu+1}
	\simeq
	-\frac{\gamma B_{c}}{2\nu}
	\left(
		\frac{\eta_{+}^{2\nu}}{S_{\nu-1}}
		+
		\frac{2\nu}{\nu+1}\eta_{+}^{\nu+1}
	\right).\\
	&\simeq
	-\frac{\gamma B_{c}}{2}
	\Theta(t)
	\begin{cases}
		\nu S_{\nu-1}t^{2},
		&0<\nu<1,\\[4pt]
		\dfrac{S_{0}}{S_{0}+1}t^{2},
		&\nu=1,\\[8pt]
		\dfrac{2}{\nu+1}t^{\nu+1},
		&\nu>1.
	\end{cases}
\end{equation*}

进而磁化强度、磁化率的奇异项为
\begin{equation*}
	M_{\mathrm{sing}}&
	=
	-\frac{\partial f_{\mathrm{sing}}}{\partial B}
	\simeq
	\frac{\nu+1}{\nu B_{c}}\frac{S_{\nu}}{S_{\nu-1}}
	\frac{\partial f_{\mathrm{sing}}}{\partial t}
	\simeq
	-\gamma\Theta(t)
	\begin{cases}
		(\nu+1)S_{\nu}t,
		&0<\nu<1,\\[4pt]
		\dfrac{2S_{1}}{S_{0}+1}t,
		&\nu=1,\\[8pt]
		\dfrac{(\nu+1)S_{\nu}}{\nu S_{\nu-1}}t^{\nu},
		&\nu>1,
	\end{cases}\\[8pt]
	\chi_{\mathrm{sing}}&
	=
	-\frac{\partial^{2}f_{\mathrm{sing}}}{\partial B^{2}}
	\simeq
	-\left(
	\frac{\nu+1}{\nu B_{c}}\frac{S_{\nu}}{S_{\nu-1}}
	\right)^{2}
	\frac{\partial^{2}f_{\mathrm{sing}}}{\partial t^{2}}\simeq
	\frac{\gamma}{B_{c}}
	\Theta(t)
	\begin{cases}
		\dfrac{(\nu+1)^{2}S_{\nu}^{2}}{\nu S_{\nu-1}},
		&0<\nu<1,\\[12pt]
		\dfrac{4S_{1}^{2}}{S_{0}(S_{0}+1)},
		&\nu=1,\\[12pt]
		\dfrac{(\nu+1)^{2}S_{\nu}^{2}}{\nu S_{\nu-1}^{2}}t^{\nu-1},
		&\nu>1.
	\end{cases}
\end{equation*}

因此，当\(0<\nu\le1\)，即\(2<d\le4\)时，磁化率出现有限跃变；当\(\nu>1\)，即\(d>4\)时，磁化率连续，非解析性出现在更高阶导数，且可能出现发散项.
